\documentclass[10pt,a4paper]{amsart}
\usepackage[margin=19mm]{geometry}
\usepackage{amsmath,amssymb,mathtools}
\usepackage[scr=rsfs,cal=euler]{mathalfa}
\usepackage{xfrac}
\usepackage[unicode,hidelinks,bookmarks=false]{hyperref}
\newcommand{\ie}{i.e.\ }
\newcommand{\cf}{cf.\ }
\newcommand{\resp}{respectively\ }
\newcommand{\Subsection}{\subsection}
\providecommand{\nobreakdash}{\nobreak\hbox{-}}
\setlength{\emergencystretch}{2em}
\multlinegap0pt
\newtheorem{theorem}{Theorem}
\newtheorem{lemma}[theorem]{Lemma}
\newtheorem{proposition}[theorem]{Proposition}
\newtheorem{definition}[theorem]{Definition}
\newcommand{\ur}{\mathrm{un}}
\newcommand{\Q}{\mathbb{Q}}
\newcommand{\Z}{\mathbb{Z}}
\newcommand{\fS}{\mathfrak{S}}
\newcommand{\bJ}{\mathbf{J}}
\newcommand{\bA}{\mathbf{A}}
\newcommand{\Sel}{\operatorname{Sel}}
\renewcommand{\O}{\mathcal{O}}
\title[Theorem 4.14]{On a parity result for the symmetric square of modular forms with congruent residual representations: the imprimitive signed $\lambda$-invariant congruence (Theorem 4.14)}
\author{Jishnu Ray}
\date{Source: arXiv:2406.03050v3 (CC BY 4.0)}
\begin{document}
\maketitle

\noindent\emph{Source and licence.} On a parity result for the symmetric square of modular forms with congruent residual representations, by Jishnu Ray. Version arXiv:2406.03050v3
(CC BY 4.0), distributed by arXiv under the Creative Commons Attribution 4.0 International licence,
http://creativecommons.org/licenses/by/4.0/, as recorded in the arXiv licence metadata for this exact
version and shown on the abstract page https://arxiv.org/abs/2406.03050v3. Full licence notice:
\texttt{licenses/arxiv-2406.03050-CC-BY-4.0.txt}.\par\smallskip

\noindent\textit{Editorial scope.} Counted unit: the article's Theorem 4.14 (source0.tex lines 977--980, source label \texttt{prop: lambda}), the non-ordinary half of the paper's parity theorem: for two modular forms with isomorphic residual representations and $a_p=0$, the $\mu$-invariant of the signed Selmer group of $\bA_{f_1,\psi}$ vanishes if and only if that of $\bA_{f_2,\psi}$ vanishes, and when both vanish the imprimitive signed $\lambda$-invariants coincide. The article's complete original proof of Theorem 4.14 is delivered with it (source0.tex lines 981--985, one \texttt{proof} environment of five lines, adjacent to the statement, with no gap and no deferral). No mathematics is written, repaired or re-derived: the statement and the proof are the article's own lines, in the article's own notation, and no retained line is substituted, re-flowed or omitted, so nothing is removed and no omission record is needed (the printed wording of line 984, ``one case easily show'', is kept literally). Retained uncounted as the source-local context the proof needs: the paragraph deriving the $[\mathfrak{P}]$-torsion identifications (lines 916--936, one contiguous run) and the statement of the residual congruence of the two imprimitive signed Selmer groups (lines 948--951, printed as Proposition 4.13, without its proof). Inside that run are the article's definition of the $\Sigma_0$-imprimitive signed Selmer group (lines 924--928, printed as Definition 4.10), the isomorphism it induces on $[\mathfrak{P}]$-torsion (lines 929--932, printed as equation (19), source label \texttt{eq:jishnu18}) and the statement of the cotorsion lemma the proof invokes (lines 933--936, printed as Lemma 4.11, source label \texttt{lem:4.9}). Every cross-reference of every retained line resolves inside this document: \texttt{\textbackslash ref\{def:imprimitive\}}, \texttt{\textbackslash eqref\{eq:jishnu18\}}, \texttt{\textbackslash ref\{lem:4.9\}} and \texttt{\textbackslash ref\{prop:congprop\}}. No retained line contains a citation, so the wrapper carries no bibliography and no bibliography entry. No retained line contains a figure environment, a graphics-inclusion command or drawing code, so no image is delivered and the package declares no asset dependency.

Editorial formatting only, in addition: an AMS \texttt{amsart} preamble that restates the article's own declarations and macros used by the retained lines -- the theorem-like environments \texttt{theorem}, \texttt{lemma}, \texttt{proposition} and \texttt{definition} sharing one counter as in the article, and the macros \texttt{\textbackslash Q}, \texttt{\textbackslash Z}, \texttt{\textbackslash fS}, \texttt{\textbackslash bA}, \texttt{\textbackslash bJ}, \texttt{\textbackslash Sel}, \texttt{\textbackslash O} and \texttt{\textbackslash ur} copied from the article's own preamble (source0.tex lines 65, 89, 92, 116, 121, 206, 312 and 314) -- and consecutive numbering of the retained environments in their order of appearance, so that in this document the retained Definition prints as Definition 1, the retained displayed isomorphism as (1), the retained Lemma as Lemma 2, the retained Proposition as Proposition 3 and the counted Theorem as Theorem 4, whereas the article numbers its theorem-like environments by section and prints the same five items as Definition 4.10, (19), Lemma 4.11, Proposition 4.13 and Theorem 4.14. No section heading is delivered, so the wrapper numbers the retained environments consecutively instead of per section.

The complete native source of the article as distributed in the arXiv e-print for this exact version is bundled unchanged as \texttt{sources/160309/source0.tex}: it is byte identical to the e-print file \texttt{revisionIJNT2.tex} with the article's own macro file \texttt{canonicalheader.tex} spliced in where the article reads it, except that the pipeline's mechanical concatenation prepends one header comment line naming the spliced files. The article is an arXiv preprint typeset with the AMS \texttt{amsart} class; no publisher class or style file exists for it and none is redistributed or used.
From the following exact sequence 
$$0 \rightarrow \bA_{f, \psi}[\mathfrak{P}] \rightarrow \bA_{f, \psi} \rightarrow \bA_{f, \psi} \rightarrow 0 $$
we obtain that the sequence 
$$ 0 \rightarrow H^0(\Q_\infty, \bA_{f, \psi})/\mathfrak{P} \rightarrow H^1(\Q_\infty, \bA_{f, \psi}[\mathfrak{P}] ) \rightarrow H^1(\Q_\infty,\bA_{f, \psi})[\mathfrak{P}] \rightarrow 0. $$
As $H^0(\Q_\infty, \bA_{f, \psi})=0$ by \textbf{(invS)}, we obtain $$H^1(\Q_\infty, \bA_{f, \psi}[\mathfrak{P}] ) \cong H^1(\Q_\infty,\bA_{f, \psi})[\mathfrak{P}]. $$ The same proof also gives that, for a prime $v \mid p$, 
$$H^1((\Q_\infty)_v, \bA_{f, \psi}[\mathfrak{P}] ) \cong H^1((\Q_\infty)_v,\bA_{f, \psi})[\mathfrak{P}]. $$
Define $$H^1_{\fS}((\Q_\infty)_v, \bA_{f, \psi}[\mathfrak{P}] ):=H^1_{\fS}((\Q_\infty)_v, \bA_{f, \psi})[\mathfrak{P}].$$
It is also well-known that $H^1_{\ur}((\Q_\infty)_v,\bA_{f, \psi}[\mathfrak{P}])=H^1_{\ur}((\Q_\infty)_v,\bA_{f, \psi})[\mathfrak{P}]$ for primes $v\nmid p$ and $v\not\in \Sigma.$ 
\begin{definition}\label{def:imprimitive} The $\Sigma_0$-imprimitive signed Selmer group $\Sel_{\fS}^{\Sigma_0}(\bA_{f,\psi}/\Q_\infty)$ for the Galois module $\bA_{f,\psi}$ is define as the kernel of the following map 
$$H^1(\Q_\infty,\bA_{f, \psi}) \rightarrow \prod_{v \mid p} \frac{H^1((\Q_\infty)_v,\bA_{f, \psi})}{H_{\fS}^1((\Q_\infty)_v,\bA_{f, \psi})}  \times \prod_{v \in \Sigma\backslash \Sigma_0}\frac{H^1((\Q_\infty)_v,\bA_{f, \psi})}{H^1_{\ur}((\Q_\infty)_v,\bA_{f, \psi})}.$$
The $\Sigma_0$-imprimitive signed Selmer group $\Sel_{\fS}^{\Sigma_0}(\bA_{f,\psi}[\mathfrak{P}]/\Q_\infty)$ for the Galois module $\bA_{f,\psi}[\mathfrak{P}]$ is define as the kernel of the following map 
$$H^1(\Q_\infty,\bA_{f, \psi}[\mathfrak{P}]) \rightarrow \prod_{v \mid p} \frac{H^1((\Q_\infty)_v,\bA_{f, \psi}[\mathfrak{P}])}{H_{\fS}^1((\Q_\infty)_v,\bA_{f, \psi}[\mathfrak{P}])}  \times \prod_{v \in \Sigma\backslash \Sigma_0}\frac{H^1((\Q_\infty)_v,\bA_{f, \psi}[\mathfrak{P}])}{H^1_{\ur}((\Q_\infty)_v,\bA_{f, \psi}[\mathfrak{P}])}.$$  
\end{definition}
By the discussion before definition \ref{def:imprimitive}, it follows that the local conditions defining the Selmer groups $\Sel_{\fS}^{\Sigma_0}(\bA_{f,\psi}/\Q_\infty)[\mathfrak{P}]$ and $ \Sel_{\fS}^{\Sigma_0}(\bA_{f,\psi}[\mathfrak{P}]/\Q_\infty)$ are the same. Hence there is a $\Lambda$-module isomorphism 
\begin{equation}\label{eq:jishnu18}
    \Sel_{\fS}^{\Sigma_0}(\bA_{f,\psi}/\Q_\infty)[\mathfrak{P}] \cong  \Sel_{\fS}^{\Sigma_0}(\bA_{f,\psi}[\mathfrak{P}]/\Q_\infty).
\end{equation}
\begin{lemma}\label{lem:4.9}
    Assume that $\Sel_{\fS}(\bA_{f, \psi})$ and $\Sel_{\fS}(\bJ_{f, \psi^{-1},-t})$ are cotorsion $\Lambda$-modules for some fixed $t\in \Z$. Then the Selmer group $\Sel^{\Sigma_0}_{\fS}(\bA_{f, \psi})$ is also $\Lambda$-cotorsion and 
     $$\mu(\Sel_{\fS}(\bA_{f, \psi})) =\mu(\Sel^{\Sigma_0}_{\fS}(\bA_{f, \psi})).$$
\end{lemma}

\begin{proposition}\label{prop:congprop}
    We have the following isomorphism as $\Lambda$-modules.
    $$\Sel_{\fS}^{\Sigma_0}(\bA_{f_1, \psi}[\mathfrak{P}]/\Q_\infty)\cong \Sel_{\fS}^{\Sigma_0}(\bA_{f_2, \psi}[\mathfrak{P}]/\Q_\infty).$$
\end{proposition}

\begin{theorem}\label{prop: lambda}
Assume that the hypothesis of Lemma \ref{lem:4.9} is true for both $f_1$ and $f_2$.  Then the $\mu$-invariant of $\Sel_{\fS}(\bA_{f_1, \psi}/\Q_\infty)$ vanishes if and only if the $\mu$-invariant of $\Sel_{\fS}(\bA_{f_2, \psi}/\Q_\infty)$ vanishes. When these $\mu$-invariants are trivial, then the imprimitive signed $\lambda$-invariants coincide, i.e. 
$$\lambda (\Sel^{\Sigma_0}_{\fS}(\bA_{f_1, \psi}/\Q_\infty))=\lambda (\Sel^{\Sigma_0}_{\fS}(\bA_{f_2, \psi}/\Q_\infty)).$$
\end{theorem}

\begin{proof}
    By Lemma \ref{lem:4.9}, the $\mu$-invariant of $\Sel_{\fS}(\bA_{f_i, \psi}/\Q_\infty)$ vanishes if and only if $\Sel_{\fS}^{\Sigma_0}(\bA_{f_i, \psi}/\Q_\infty)$ is cofinitely generated as an $\O$-module which is true if and only if $\Sel_{\fS}^{\Sigma_0}(\bA_{f_i, \psi}/\Q_\infty)[\mathfrak{P}]$ is finite. By \eqref{eq:jishnu18}, $\Sel_{\fS}^{\Sigma_0}(\bA_{f_i, \psi}/\Q_\infty)[\mathfrak{P}]$ is finite if and only if  $\Sel_{\fS}^{\Sigma_0}(\bA_{f_i, \psi}[\mathfrak{P}]/\Q_\infty)$ is finite.  Then Proposition \ref{prop:congprop} gives that $\Sel_{\fS}^{\Sigma_0}(\bA_{f_1, \psi}[\mathfrak{P}]/\Q_\infty)$ is finite if and only if $\Sel_{\fS}^{\Sigma_0}(\bA_{f_2, \psi}[\mathfrak{P}]/\Q_\infty)$ is finite. 
\vspace{.2cm}
Since the imprimitive Selmer group $\Sel^{\Sigma_0}_{\fS}(\bA_{f_i, \psi}/\Q_\infty)$ has no proper $\Lambda$-submodule of finite index, one case easily show that $$\lambda(\Sel^{\Sigma_0}_{\fS}(\bA_{f_i, \psi}/\Q_\infty))=\dim_{\mathbb{F}_p}\Sel_{\fS}^{\Sigma_0}(\bA_{f_i, \psi}/\Q_\infty)[\mathfrak{P}].$$ But $\Sel_{\fS}^{\Sigma_0}(\bA_{f_1, \psi}/\Q_\infty)[\mathfrak{P}]\cong \Sel_{\fS}^{\Sigma_0}(\bA_{f_2, \psi}/\Q_\infty)[\mathfrak{P}]$ and hence their $\mathbb{F}_p$-dimensions coincide. 
\end{proof}

\end{document}
